\subsection{(\S4) Archetypal aspects of masculine adaptation}

\subsubsection{Prologue: masculine psychology}

\begin{enumerate}

\item \fbox{\it module and role} The masculine mind, as I have encountered it
in men who have submitted their mental processes to me for therapeutic
analysis, organizes itself according to separate modules of identity, each for
a time lived as if it were his only self—that is, identified with, defended,
and reified into a role by which the man who latches onto it tries to
live—unless of course the man is preferring to dissolve his mind in a Dionysian
way with the aid of alcohol, drugs, music, a sex partner, or even a mood, in
which case his femininity emerges to release the disorganized flow that Woolf
describes, though often without her accompanying ambition to find a self that
will make the flux cohere.

\item \fbox{\it masculine and Jung's psychology} For this reason, the
categorical nature of Jung's psychology of the unconscious has often been one
of its most compelling features for men. In the archetypal complexes that Jung
was the first to identify as ingredients of personality, a man can recognize
his part personalities—the roles and values that he might choose to take up
when trying to adapt to living in the world in a masculine way. Men often focus
their efforts in therapy on living up to one particular role. Part of the job
of the psychotherapist in that situation is to see how the patient's psyche is
responding to his unconscious decision to disregard the other identities and
values that are available to him.

\item \fbox{\it reader suggestion} As the reader explores the different
archetypes that I present here, she or he should keep in mind that the role
each implies could seem to some men like just the way a man is supposed to be.

\end{enumerate}

\subsubsection{The hero}

\begin{enumerate}

    \item \fbox{\it definition} We might start, as the study of archetypes
        usually does, with the hero. “The hero,” Jung (1989) tells us, “is the
        symbol of the greatest value recognized by us” (p. 57). In our culture,
        with its patriarchal system of education modeled upon the Greeks and
        Romans, perhaps the greatest values imparted to men are achievement and
        responsibility. These standards cannot be met by someone who lacks
        autonomy, accuracy, and self-control—aspects of self-reliant agency
        that are exemplified in the myths of the hero (Segal, 1990). Stories of
        the hero continue to have wide appeal in our culture and to exert
        profound psychological influence. From the standpoint of depth
        psychology, the hero is the archetypal role model for ego-identity,
        described for his fellow psychoanalysts by Erik Erikson as “the accrued
        experience of the ego’s ability to integrate all identifications with
        the vicissitudes of the libido, with the aptitudes developed out of
        endowment, and with the opportunities offered in social roles”
        (Erikson, 1963, p. 261).

    \item \fbox{\it hero development example} What is significant for the
        psychology of a man’s development as we encounter it in analytical work
        is how often the image of the hero will appear in the psyche of a
        client beginning to consolidate his identity out of the lost state of
        mind that has caused him to seek treatment. One of my clients first
        came to me when he was 19; he was recovering from a brief psychotic
        breakdown in college. His need to develop his own psyche separately
        from the excessive superego expectations of his parents soon became a
        theme of the therapy. After the first few months of our work, he
        dreamed he was at a political rally where a new leader was chosen by
        the vote of those present. After he had this dream, his life changed
        and it was almost impossible to believe that the patient had ever been
        psychotic. For years thereafter he consistently displayed a strong,
        decisively masculine ego, and soon moved forward into high levels of
        achievement, marriage, and a successful professional career. The hero
        figure, elected by the political party in his dream, evidently went on
        to take office in his psyche: in retrospect, it was the image of the
        emergence of ego strength in this patient.

    \item \fbox{\it characteristics} Jungian theory offers a way of
        understanding the hero figure in a man’s psyche that is even more
        precise: the hero symbolizes what Jung calls the dominant function of
        consciousness and is thus an image of the cutting edge of the ego. For
        Jung (1939/1959), “Consciousness needs a centre, an ego to which
        something is conscious” (¶506). But this is not all. Will is perhaps
        the most characteristic attribute of the ego once it assumes its place
        at the center of consciousness, and I could certainly see that strong
        will emerge in my patient, who effectively took charge of his therapy
        after having this dream, by asking me to hurry up, be more concise, get
        to the point, and so on.

    \item \fbox{\it liberation via hero development} His ability to take heroic
        precedence over me in directing the therapy permitted him to develop a
        dominant ego that would free him from his parents and allow him to
        become his own man.

\end{enumerate}

\subsubsection{The puer aeternus}

\begin{enumerate}

\item \fbox{\it definition} In our time, the term puer aeternus has staged a
return in Jungian psychology to indicate a man who “remains too long in
adolescent psychology” (von Franz, 1970, p. 1) and, although charming and
promising, is ultimately unreliable.

\item \fbox{\it client example} He looked like a young Lord Byron, and his air,
in the outpatient clinic where we met for psychotherapy, was imperious and
dismissive. (...) Since he is special, he sees “no need to adapt.” He evinces
an: arrogant attitude towards other people due to both an inferiority complex
and false feelings of superiority. (...) In relationships, too, the right
partner is never found, because no one can match the ideal person the patient
is looking for.

\item \fbox{\it how to help this client} One way to help such a person is to
engage his feelings about the therapeutic relationship. In my patient’s case,
getting him to express his anger at having to be in therapy at all helped to
ground the work and move it in the direction of real adaptation. It turned out
that he was irritated with the power imbalance of therapy and that this was why
he was treating me with such disdain. Our work finally addressed this reality
when I insisted that he show a measure of respect for my role as therapist.
Predictably, this made him furious, but his fury brought his human ego into the
relationship. In Jungian terms, the centerpiece of his ego was introverted
feeling, which is very sensitive to imbalances of power. He told me in no
uncertain terms that it was I who was being pompous. It would have been
tempting to put this down as projective identification, but accepting his
resentment of my assertion of power was the more helpful course.

\item \fbox{\it changes seen afterwards} Discriminating the appropriate and
inappropriate uses of power is something people with strong introverted feeling
are good at; this function is often symbolized in dreams as the judge. As the
patient became less ‘above it all’ and more directly confrontational with me,
he gradually began to make more realistic judgments as to what he needed to do
in his own educational, vocational, and romantic life.

\item \fbox{\it deeper empathy} This bit of case illustrates how first
recognizing the archetype effectively possessing the person and then
identifying the psychological type of the patient’s real ego, which has been
buried in the archetypal presentation, can be ways for the Jungian therapist to
more deeply empathize with the nature of a patient’s struggle in life.

\end{enumerate}

\subsubsection{The trickster}

\begin{enumerate}

\item \fbox{\it definition} A third archetype that not infrequently makes its
appearance in the psychotherapist’s consulting room is the trickster. It is one
that is problematic for analytical work, because the very ethos of this kind of
psychotherapy is a commitment to being sincere and vulnerable. A man identified
with the trickster archetype is anything but sincere and vulnerable. The
archetypal thrust of the trickster’s restless energy is to resist having to do
things the way others think is right. The trickster toys with expectations,
flouting rules that attempt to uphold standards of behavior and questioning the
values of those who defend the standards (Radin, 1972).

\item \fbox{\it double bind example} It was this communication that impressed
him and motivated him to change his life, and for more reasons than I could
have known at the time. A year and a half into his therapy, he mentioned that
when, in high school, he had told his stepfather that he dreamed some day of
going to graduate school, the stepfather had coldly replied, "I don't care what
you do; I don't care if you sell your body." In the therapy, the young man had
finally met someone who didn't accept the terms of the double bind he had been
placed in by the stepfather's remark, that the only way for him to get an
education, including even a psychotherapy, was to disregard his own feelings.

\item \fbox{\it for therapist} When the trickster appears out of the psyche of
a patient, its calling card is usually a display of the archetype’s capacity to
put both others and oneself in a double bind (Beebe, 1981, p. 38). The patron
divinity of alchemy, the wily Roman god Mercurius, is often referred to as
“duplex and his main characteristic is duplicity” (Jung, 1948/1967, p. 217).
Contemporary men who have aligned their survival strategies with the trickster
archetype can be quite duplicitous when urging their psychotherapists to
support personal choices that are, in fact, highly untherapeutic. The
analytical therapist will want to recognize that the trickster is operating in
the patient, but that is not enough. The therapist will have to have integrated
enough trickster of his or her own to be able to turn the double bind around
and reverse the terms of the hard bargain that the trickster in the patient is
trying to drive. When the patient initially told me that he could not pay for
the therapy unless he went on working as a prostitute, I felt the cage of the
double bind closing in around me. By making him see that he would not be able
to have me as a therapist if he continued his prostitution, I managed to turn
the trap back on him and by so doing signaled him as well that I was at least
potentially equal to the challenge of being his analyst.

\end{enumerate}

\subsubsection{The shadow}

TODO
